\documentclass[../../main.tex]{subfiles}
\input{0-theme/default}

\begin{document}

% ==========================
\section{Pyrolysis}
% ==========================

%%%%%%%%%%%%%%%%%%%%%%%%%%%%%%%%%%%%%%%%%%%%%%%%%%%%%%%%%%%%%%%%%
\subsection{Introducción}
\begin{frame}{Introducción}
\footnotesize
\setlength{\parskip}{2pt}

La \textbf{pirólisis} es la descomposición termoquímica de la biomasa que genera una variedad de productos útiles, ya sea en ausencia total de agentes oxidantes o con un suministro tan limitado que no permite la gasificación en forma apreciable. Además, constituye varias etapas iniciales del proceso de gasificación.

Durante la pirólisis, las grandes moléculas hidrocarbonadas complejas de la biomasa se descomponen en moléculas más pequeñas y sencillas de gas, líquido y carbón.

Procesos relacionados o que se solapan: \emph{craqueo, desvolatilización, carbonización, torrefacción, destilación seca, destilación destructiva y termólisis}. No se asemeja a la \textbf{gasificación}, que implica reacciones con un agente externo (medio de gasificación).

\textbf{Rangos de temperatura típicos:}
\begin{itemize}
\item Pirólisis: 300--650~°C
\item Gasificación: 800--1000~°C
\item Torrefacción: 200--300~°C
\end{itemize}

El \textbf{biocarbón} es el producto sólido resultante de la pirólisis; ofrece una alternativa a la quema completa de residuos, pues puede retener carbono en el suelo durante cientos de años.

\end{frame}

\begin{frame}{Antecedentes Históricos}
\begin{itemize}
\item \textbf{Era preindustrial:} El carbón vegetal (pirólisis lenta de madera) fue esencial para extraer hierro a partir del mineral.
\item Se utilizaban hornos tipo \emph{colmena} para producir carbón de biomasa.
\item Cuando la madera escaseó, el carbón mineral de minas subterráneas sustituyó al carbón vegetal.
\item \textbf{Innovación de Abraham Gesner (Halifax, 1840s):}
\begin{itemize}
\item Destiló trozos de carbón a \SI{427}{\celsius}, purificó con \(H_2SO_4\) y cal, y redestiló.
\item Obtuvo un líquido claro que quemaba con llama brillante: lo llamó \textbf{queroseno}.
\end{itemize}
\item El queroseno sustituyó el aceite de ballena, reduciendo la flota ballenera (>730 barcos en 1846) y evitando la extinción.
\item Fue el primer combustible líquido transportable, revolucionando la iluminación mundial.
\end{itemize}
\end{frame}
\begin{frame}{Pyrolysis}
\small
\begin{itemize}
  \item \textbf{Definición:} Calentamiento rápido de biomasa en ausencia de O\textsubscript{2} hasta la \emph{temperatura de pirolisis}, manteniéndola para obtener gases no condensables, carbón (\emph{char}) y un producto líquido (bio‑aceite).
  \item El bio‑aceite es el producto de mayor interés; su naturaleza depende de la temperatura y la velocidad de calentamiento.
  \item Productos iniciales: vapores condensables + \emph{char}. Los vapores pueden craquearse en CO, CO\textsubscript{2}, H\textsubscript{2}, CH\textsubscript{4}, más líquido y \emph{char} adicional.
  \item Mecanismos: reacciones homogéneas en fase gas y heterogéneas gas‑sólido.
  \item Reacción genérica:
    \[
      \mathrm{C_{n}H_{m}O_{p}}(biomasa)  \xrightarrow{\text{calor}} 
      \sum\mathrm{C_{x}H_{y}O_{z}}_{\text{(líq)}} +
     \sum \mathrm{C_{a}H_{b}O_{c}}_{\text{(gas)}} + \mathrm{H_2O} + \mathrm{C}_{\text{(char)}}
    \]
  \item La pyrolysis es un paso previo rápido y esencial dentro de un gasificador.
 
\end{itemize}
\end{frame}
\begin{frame}{planta de pyrólysis}
 \textbf{Esquema de lecho fluidizado:} Biomasa alimentada sobre sólidos calientes; vapores salen y se enfrían; el bio‑aceite condensa; los gases no condensables pueden reciclarse; el \emph{char} puede recuperarse o quemarse para generar calor.
    \begin{figure}
        \centering
        \includegraphics[width=0.55\linewidth]{Energias Alternativas/3-Biomass/img/simple_planta_pyrolysis.png}
        \small
        \caption{Esquema simplificado de una planta de pyrólysis.
}
        \label{fig:enter-label}
    \end{figure}
\end{frame}



\begin{frame}{Productos de la Pirólisis}
\scriptsize
\begin{itemize}
  \item La pirólisis descompone grandes moléculas complejas en moléculas más pequeñas.
  \item \textbf{Productos principales:}
  \begin{enumerate}
    \item \textbf{Líquido}: alquitranes, hidrocarburos pesados y agua.
    \item \textbf{Sólido}: char (carbón sólido).
    \item \textbf{Gas}:
    \begin{itemize}
      \item \textit{CO\textsubscript{2}}: dióxido de carbono
      \item \textit{H\textsubscript{2}O}: vapor de agua
      \item \textit{CO}: monóxido de carbono
      \item \textit{C\textsubscript{2}H\textsubscript{2}}: acetileno
      \item \textit{C\textsubscript{2}H\textsubscript{4}}: etileno
      \item \textit{C\textsubscript{2}H\textsubscript{6}}: etano
      \item \textit{C\textsubscript{6}H\textsubscript{6}}: benceno
    \end{itemize}
  \end{enumerate}
  \item Las fracciones relativas dependen de la velocidad de calentamiento y la temperatura final (\SI{300}{\degreeCelsius}–\SI{650}{\degreeCelsius}).
  \item No confundir con la \emph{materia volátil}: en el análisis próximo, líquidos y gases se agrupan como materia volátil, y el char como carbono fijo.
  \item Normas ASTM (p. ej. \textbf{D‑3172}, \textbf{D‑3175}) definen condiciones estándar: horno a \SI{950}{\degreeCelsius} durante \textbf{7 min} para medir materia volátil.
\end{itemize}
\end{frame}

\begin{frame}{Fracción líquida de la pyrolysis}
\begin{itemize}
\item Fluido negro y viscoso denominado \emph{alquitrán}, \emph{bio‑oil} o \emph{biocrudo}.
\item Contenido de agua: \SI{\approx 20}{\percent}.
\item Compuesto principalmente por compuestos fenólicos homólogos.
\item Poder calorífico inferior (LHV):
\begin{itemize}
\item Biomasa original (base seca): \SIrange{19.5}{21}{MJ,kg^{-1}}.
\item Bio‑oil (base húmeda): \SIrange{13}{18}{MJ,kg^{-1}}.
\end{itemize}
\item \textbf{Formación}: despolimerización y fragmentación rápidas de celulosa, hemicelulosa y lignina, seguida de enfriamiento instantáneo (\emph{quenching}) para "congelar" los productos intermedios.
\item \textbf{Naturaleza}: microemulsión
\begin{itemize}
\item Fase continua: solución acuosa de productos de celulosa y hemicelulosa + moléculas pequeñas de lignina.
\item Fase discontinua: macromoléculas de lignina pirrolítica.
\end{itemize}
\end{itemize}
\end{frame}

\begin{frame}{Fracción líquida de la pyrolysis}
\begin{itemize}
     \item Peso molecular del bio‑oil condensado: $>,500$~Da.
\item \textbf{Principales familias de compuestos} (Piskorz \emph{et al.}, 1988):
\end{itemize}
\begin{itemize}
\item Hidroxi‑aldehídos
\item Hidroxi‑cetonas
\item Azúcares y deshidro‑azúcares
\item Ácidos carboxílicos
\item Compuestos fenólicos
\end{itemize}
\end{frame}

%------------------------------
\begin{frame}{Fracción Sólida: Biocarbón}
\begin{itemize}
\item Producto sólido resultante de la \textbf{pyrolysis}.
\item Composición aproximada: \textbf{85,\% C} con trazas de O y H; mínima ceniza inorgánica.
\item \textbf{Poder calorífico inferior (LHV)}: \SI{32}{\mega\joule\per\kilogram} (Diebold y Bridgwater, 1997), superior al de la biomasa original y al bio‑oil.
\item Presenta una \textbf{gran área superficial porosa}, útil para adsorción y almacenamiento de carbono.
\end{itemize}
\end{frame}
%----------------------------------------
\begin{frame}{Gas en la Pyrolysis}
\small
\begin{itemize}
\item \textbf{Descomposición primaria} genera:
\begin{itemize}
\item \textbf{Vapores condensables}: moléculas pesadas que se licuan al enfriarse y aumentan el rendimiento líquido.
\item \textbf{Gases no condensables (primarios)}: \ce{CO2}, \ce{CO}, \ce{CH4}, \ce{C2H6}, \ce{C2H4}.
\end{itemize}
\item El \textit{craqueo secundario} del vapor a mayor $T$ produce \textbf{gases secundarios} adicionales.
\item Producto final = mezcla de gases primarios + secundarios.
\item $\text{LHV}_{\text{primario}} \approx \SI{11}{MJ.Nm^{-3}}$; tras craqueo severo, $\text{LHV}_{\text{pyrolysis}} \approx \SI{20}{MJ.Nm^{-3}}$ (\textit{Diebold \& Bridgwater}, 1997).
\end{itemize}
\end{frame}


%-----------------------------------------------------------------
\begin{frame}{Clasificación de la Pyrolysis}
\begin{block}{Según la velocidad de calentamiento}
\begin{itemize}
\item \textbf{Lenta:} $t_{\text{heating}} \gg t_{r}$ – orientada a producción de char.
\begin{itemize}
\item \emph{Torrefacción}: \SI{200}{--,300},\si{\degreeCelsius}
\item \emph{Carbonización}: $>,\SI{400}{\degreeCelsius}$
\end{itemize}
\item \textbf{Rápida:} $t_{\text{heating}} \ll t_{r}$ – orientada a bio‑oil y gas.
\begin{itemize}
\item \emph{Flash}: residencia $<,\SI{1}{s}$
\item \emph{Ultrarapid}: residencia $<,\SI{0.5}{s}$
\end{itemize}
\end{itemize}
\end{block}
\vspace{0.5em}
\begin{block}{Según el medio no oxidante}
\begin{itemize}
\item \emph{Hydrous}: en H\textsubscript{2}O – producción de químicos.
\item \emph{Hydro‑}: en H\textsubscript{2} – bio‑oil.
\item \emph{Methano‑}: en CH\textsubscript{4} – obtención de químicos.
\end{itemize}
\end{block}
\end{frame}

%-----------------------------------------------------------------
\begin{frame}{Tabla – Características de procesos térmicos}
\small
\begin{tabular}{@{}lcccc@{}}
\toprule
\textbf{Proceso} & \textbf{Residencia} & \textbf{dT/dt} & \textbf{T final (°C)} & \textbf{Productos} \\
\midrule
Torrefacción   & 10--60 min & Muy baja  & 280      & Biomasa torrefactada \\
Carbonización  & Días       & Muy baja  & $>400$   & Carbón vegetal \\
Rápida         & $<2$ s     & Muy alta  & $\sim500$ & Bio‑oil \\
Flash          & $<1$ s     & Alta      & $<650$   & Bio‑oil, químicos, gas \\
Ultrarapid     & $<0.5$ s   & Muy alta  & $\sim1000$ & Químicos, gas \\
Vacío          & 2--30 s    & Media     & 400      & Bio‑oil \\
Hydro‑         & $<10$ s    & Alta      & $<500$   & Bio‑oil \\
Methano‑       & $<10$ s    & Alta      & $>700$   & Químicos \\
\bottomrule
\end{tabular}
\end{frame}

\begin{frame}{Pirólisis Lenta y Convencional}
\footnotesize
\textbf{Pirólisis lenta (Carbonización):}
\begin{itemize}
  \item Proceso más antiguo de pirólisis, utilizado desde hace miles de años.
  \item Objetivo principal: producción de \textbf{char} o carbón vegetal.
  \item La biomasa se calienta lentamente en ausencia de oxígeno a $\sim\SI{400}{\degreeCelsius}$ durante un periodo prolongado (incluso días).
  \item En hornos tipo colmena, los troncos se apilan y cubren con arcilla; se permite entrada limitada de oxígeno para combustión parcial.
  \item El calor es generado por un pequeño fuego interno; el sistema permanece bien aislado.
  \item El tiempo prolongado favorece la conversión de vapores condensables en char y gases no condensables.
\end{itemize}

\textbf{Pirólisis convencional:}
\begin{itemize}
  \item Produce \textbf{gas, líquido y char} como productos principales.
  \item Calentamiento moderado hasta $\sim\SI{600}{\degreeCelsius}$.
  \item Tiempo de residencia del producto: del orden de minutos.
\end{itemize}
\end{frame}

\begin{frame}{Pirólisis Rápida}
\footnotesize
\begin{itemize}
  \item Objetivo principal: \textbf{maximizar la producción de bioaceite (líquido)}.
  \item La biomasa se calienta rápidamente hasta alcanzar la temperatura pico antes de descomponerse.
  \item \textbf{Velocidad de calentamiento:} entre \SI{1000}{\degreeCelsius\per\second} y \SI{10000}{\degreeCelsius\per\second}.
  \item \textbf{Temperatura pico:}
  \begin{itemize}
    \item $\leq \SI{650}{\degreeCelsius}$ para producir bioaceite.
    \item Hasta \SI{1000}{\degreeCelsius} si se desea producir gas.
  \end{itemize}
  \item Se emplean comúnmente \textbf{lechos fluidizados}.
\end{itemize}

\textbf{Factores clave que aumentan el rendimiento líquido:}
\begin{enumerate}
  \item Velocidad de calentamiento muy alta.
  \item Temperatura de reacción entre \SI{425}{\degreeCelsius} y \SI{600}{\degreeCelsius}.
  \item Tiempo de residencia del vapor en el reactor menor a \SI{3}{\second}.
  \item Enfriamiento rápido (quenching) del gas producto.
\end{enumerate}
\end{frame}
\begin{frame}{Pirólisis Flash}
\small
\begin{itemize}
  \item La biomasa se calienta rápidamente en ausencia de oxígeno.
  \item Rango de temperatura: \SI{450}{\degreeCelsius}–\SI{600}{\degreeCelsius}.
  \item El producto (gases condensables y no condensables) abandona el pirolizador en un \textbf{tiempo de residencia corto}: \SI{30}{ms}–\SI{1500}{ms} %\cite{Bridgwater1999}.
  \item Al enfriarse, el vapor condensable se convierte en \textbf{bioaceite (bio-oil)}.
  \item Este proceso incrementa el rendimiento de líquido y \textbf{reduce la producción de char}.
  \item Rendimiento típico de bioaceite: \textbf{70–75\,\%} del total de productos de la pirólisis.
\end{itemize}
\end{frame}
\begin{frame}{Pirólisis Ultrarrápida}
\small
\begin{itemize}
  \item Implica un \textbf{calentamiento extremadamente rápido} de la biomasa.
  \item En un método, \textbf{sólidos portadores de calor} impactan con el vapor de biomasa, generando una tasa de calentamiento muy alta (Figura 5.7C).
  \item El \textbf{enfriamiento rápido} del producto primario ocurre directamente dentro del reactor.
  \item Un \textbf{separador gas-sólido} separa los sólidos calientes de los gases no condensables y vapores del producto primario.
  \item Los sólidos se \textbf{recalientan en un combustor separado} y se transportan nuevamente al mezclador mediante un gas no oxidante.
  \item Característica clave: \textbf{tiempo de residencia corto, uniforme y controlado}.
  \item Temperatura de operación:
  \begin{itemize}
    \item \textbf{$\sim \SI{1000}{\degreeCelsius}$} para maximizar la producción de gas.
    \item \textbf{$\sim \SI{650}{\degreeCelsius}$} para maximizar la producción de líquido.
  \end{itemize}
\end{itemize}
\end{frame}
\begin{frame}{Pirólisis en Presencia de un Medio}
\footnotesize
\begin{itemize}
  \item \textbf{Hidropirólisis:}
  \begin{itemize}
    \item Descomposición térmica en atmósfera de \textbf{hidrógeno a alta presión}.
    \item Aumenta el rendimiento de volátiles y la producción de hidrocarburos de baja masa molar.
    \item Diferente a la hidrogasificación del char.
    \item El hidrógeno estabiliza radicales libres antes de que se repolimericen y formen char.
    \item Produce bioaceite con \textbf{menor contenido de oxígeno}.
  \end{itemize}

  \item \textbf{Pirólisis acuosa (hidrosa):}
  \begin{itemize}
    \item Craqueo térmico de biomasa en agua a alta temperatura.
    \item Ejemplo: conversión de desechos de pavo en hidrocarburos ligeros para combustible, fertilizantes o químicos.
    \item \textbf{Proceso en dos etapas:}
    \begin{itemize}
      \item Etapa 1: \SI{200}{\degreeCelsius}–\SI{300}{\degreeCelsius} bajo presión.
      \item Etapa 2: craqueo a \SI{500}{\degreeCelsius}.
    \end{itemize}
  \end{itemize}
\end{itemize}
\end{frame}

\begin{frame}{Rendimiento de Productos de la Pirólisis}
\footnotesize
\begin{itemize}
  \item El rendimiento de productos depende de:
  \begin{itemize}
    \item Diseño del pirolizador.
    \item Características físicas y químicas de la biomasa.
    \item Parámetros operativos:
    \begin{itemize}
      \item \textbf{Tasa de calentamiento}.
      \item \textbf{Temperatura final} (temperatura de pirólisis).
      \item \textbf{Tiempo de residencia} en la zona de reacción.
    \end{itemize}
  \end{itemize}
  
  \item Otros factores que afectan los rendimientos:
  \begin{itemize}
    \item Presión.
    \item Composición del gas ambiente.
    \item Presencia de catalizadores minerales.
  \end{itemize}

  \item \textbf{Efectos del calentamiento:}
  \begin{itemize}
    \item Calentamiento rápido → mayor cantidad de volátiles y char más reactivo.
    \item Calentamiento lento y mayor tiempo de residencia → formación de \textbf{char secundario} por reacción entre char primario y volátiles.
  \end{itemize}
\end{itemize}
\end{frame}

\begin{frame}{Influencia de la Composición de la Biomasa}
\small
\begin{itemize}
  \item La relación \textbf{H/C} de la biomasa afecta significativamente el rendimiento de la pirólisis.
  \item Los tres principales componentes de la biomasa lignocelulósica tienen rangos distintos de descomposición térmica:
  \begin{itemize}
    \item \textbf{Hemicelulosa:} \SI{150}{\degreeCelsius}–\SI{350}{\degreeCelsius}
    \item \textbf{Celulosa:} \SI{275}{\degreeCelsius}–\SI{350}{\degreeCelsius}
    \item \textbf{Lignina:} \SI{250}{\degreeCelsius}–\SI{500}{\degreeCelsius}
  \end{itemize}
  \item Cada componente se piróliza de forma distinta, con contribuciones diferentes al gas, líquido y char.
\end{itemize}
\end{frame}

\begin{frame}{Pirólisis de Celulosa y Hemicelulosa}
\small
\textbf{Celulosa:}
\begin{itemize}
  \item Fuente principal de vapores condensables.
  \item Se descompone entre \SI{300}{\degreeCelsius}–\SI{400}{\degreeCelsius}.
  \item Sin catalizadores, produce principalmente \textbf{levoglucosano} como monómero %\cite{Diebold1997}.
  \item A $T > \SI{500}{\degreeCelsius}$, el levoglucosano se volatiliza → más gas y aceite, poco char.
\end{itemize}

\vspace{0.3cm}
\textbf{Hemicelulosa:}
\begin{itemize}
  \item Menos estable, se descompone entre \SI{200}{\degreeCelsius}–\SI{300}{\degreeCelsius}.
  \item Produce más \textbf{gases no condensables} y menos alquitrán que la celulosa %\cite{Reed2002}.
\end{itemize}
\end{frame}

\begin{frame}{Pirólisis de la Lignina}
\small
\begin{itemize}
  \item Se descompone en un rango amplio: \SI{280}{\degreeCelsius}–\SI{500}{\degreeCelsius}.
  \item Tasa máxima de liberación: \SI{350}{\degreeCelsius}–\SI{450}{\degreeCelsius} %\cite{Kudo1957}.
  \item Produce más \textbf{char y compuestos aromáticos} que la celulosa %\cite{Soltes1981}.
  \item Bajo calentamiento lento a \SI{400}{\degreeCelsius}, genera hasta \textbf{40\,\% de su peso como char} %\cite{Klass1998}.
  \item También aporta al \textbf{rendimiento líquido} (aprox. 35\,\%): incluye componentes acuosos, alquitrán y \textbf{fenoles} (por ruptura de enlaces éter y C–C) %\cite{Mohan2006}.
  \item Rendimiento gaseoso: solo \textbf{10\,\%} de su peso original.
\end{itemize}
\end{frame}


\begin{frame}{Efecto de la Temperatura de Pirólisis (I)}
\small
\begin{itemize}
  \item Durante la pirólisis, la biomasa se calienta desde temperatura ambiente hasta una \textbf{temperatura máxima}, conocida como \textbf{temperatura de pirólisis}.
  \item Esta temperatura se mantiene constante hasta completar el proceso.
  \item La temperatura de pirólisis \textbf{afecta la composición y el rendimiento} de los productos generados.
  \item La \textbf{liberación de gases} durante la pirólisis de una biomasa varía significativamente con la temperatura:
  \begin{itemize}
    \item Cada gas (\ce{CO2}, \ce{CO}, \ce{H2}, \ce{CH4}) tiene un comportamiento distinto frente a cambios de temperatura.
  \end{itemize}
\end{itemize}
\end{frame}

\begin{frame}{Efecto de la Temperatura de Pirólisis (II)}
\small
\begin{itemize}
  \item La \textbf{cantidad de char} producido disminuye a medida que aumenta la temperatura de pirólisis.
  \item Ejemplo: En la pirólisis del bagazo de uva:
  \begin{itemize}
    \item Menor char a mayor temperatura.
    \item \textbf{Poder calorífico del char aumenta} con la temperatura.
    \item Esto se debe a:
    \begin{itemize}
      \item Aumento del \textbf{carbono fijo} (mayor poder calorífico).
      \item Disminución del contenido de volátiles.
    \end{itemize}
  \end{itemize}
  \item La producción de \textbf{gases no condensables} (\ce{CO2}, \ce{CO}, \ce{H2}, \ce{CH4}) también \textbf{aumenta con la temperatura}.
\end{itemize}
\end{frame}


\begin{frame}{Efecto de la Tasa de Calentamiento (I)}
\small
\begin{itemize}
  \item La \textbf{tasa de calentamiento} influye en el \textbf{rendimiento y composición} de los productos de la pirólisis.
  \item Calentamiento rápido (\SI{400}{\degreeCelsius}–\SI{600}{\degreeCelsius}):
  \begin{itemize}
    \item Mayor producción de volátiles condensables → \textbf{más líquido}.
  \end{itemize}
  \item Calentamiento lento a la misma temperatura:
  \begin{itemize}
    \item Mayor generación de \textbf{char}.
  \end{itemize}
  \item Ejemplo: pirólisis de esparto: %\cite{Debdoubi2006}:
  \begin{itemize}
    \item Tasa de calentamiento: \SI{50}{\degreeCelsius/min} → 45 \% de líquido.
    \item Tasa de calentamiento: \SI{250}{\degreeCelsius/min} → 68.5 \% de líquido.
    \item Temperatura de pirólisis: \SI{550}{\degreeCelsius}.
  \end{itemize}
  \item El \textbf{tiempo de residencia} también afecta: tiempos largos → formación de \textbf{char secundario} (reacción entre char y volátiles).
\end{itemize}
\end{frame}

\begin{frame}{Efecto de la Tasa de Calentamiento (II)}
\small
\textbf{Parámetros operativos recomendados según el producto deseado:}
\vspace{0.3cm}
\begin{enumerate}
  \item[\textbf{a.}] \textbf{Maximizar char (carbonización):}
  \begin{itemize}
    \item Tasa de calentamiento lenta: $<$ \SI{0.01}{\degreeCelsius/s}–\SI{2.0}{\degreeCelsius/s}
    \item Temperatura final baja.
    \item Tiempo de residencia del gas: largo.
  \end{itemize}

  \item[\textbf{b.}] \textbf{Maximizar líquido (pirólisis rápida):}
  \begin{itemize}
    \item Tasa de calentamiento alta.
    \item Temperatura final moderada: \SI{450}{\degreeCelsius}–\SI{600}{\degreeCelsius}.
    \item Tiempo de residencia del gas: corto.
  \end{itemize}

  \item[\textbf{c.}] \textbf{Maximizar gas:}
  \begin{itemize}
    \item Tasa de calentamiento moderada a lenta.
    \item Temperatura final alta: \SI{700}{\degreeCelsius}–\SI{900}{\degreeCelsius}.
    \item Tiempo de residencia del gas: largo.
  \end{itemize}
\end{enumerate}
\end{frame}

\begin{frame}{Efecto del Tamaño de Partícula}
\small
\begin{itemize}
  \item El \textbf{tamaño, forma, composición y estructura física} de la biomasa afectan la pirólisis mediante su influencia en la \textbf{tasa de calentamiento}.
  
  \item \textbf{Partículas finas:}
  \begin{itemize}
    \item Ofrecen menor resistencia a la salida de gases condensables.
    \item Los gases escapan antes de sufrir \textbf{craqueo secundario}.
    \item Resultado: \textbf{mayor rendimiento de líquidos}.
  \end{itemize}

  \item \textbf{Partículas grandes:}
  \begin{itemize}
    \item Mayor resistencia al escape de productos primarios.
    \item Favorecen el \textbf{craqueo secundario} → más char y gases.
    \item Usadas en métodos antiguos de producción de carbón vegetal.
  \end{itemize}
\end{itemize}
\end{frame}

\begin{frame}{Cinética de la Pirólisis}
\small
\begin{itemize}
  \item Estudiar la cinética de la pirólisis es clave para:
  \begin{itemize}
    \item El diseño ingenieril de pirolizadores o gasificadores.
    \item Comprender cómo los procesos afectan los rendimientos y la composición de los productos.
  \end{itemize}
  \item Tres procesos principales determinan la velocidad de pirólisis:
  \begin{enumerate}
    \item Cinética química.
    \item Transferencia de calor.
    \item Transferencia de masa.
  \end{enumerate}
  \item El proceso se divide en \textbf{cuatro etapas térmicas}, con límites difusos entre ellas.
\end{itemize}
\end{frame}

\begin{frame}{Etapas Físicas de la Pirólisis (I)}
\small
\textbf{1. Secado} ($\sim$\SI{100}{\degreeCelsius})
\begin{itemize}
  \item Se libera humedad libre y parte del agua ligada.
  \item El calor se conduce hacia el interior de la biomasa.
  \item En alta humedad, la lignina puede fundirse y luego solidificarse al enfriar.
  \item Usado en el curvado con vapor de madera %\cite{Diebold1997}.
\end{itemize}

\vspace{0.2cm}
\textbf{2. Etapa inicial} (\SI{100}{\degreeCelsius}–\SI{300}{\degreeCelsius})
\begin{itemize}
  \item Ocurre deshidratación exotérmica.
  \item Se liberan \ce{H2O}, \ce{CO} y \ce{CO2}.
  \item En esta etapa se lleva a cabo la \textbf{torrefacción}.
\end{itemize}
\end{frame}

\begin{frame}{Etapas Físicas de la Pirólisis (II)}
\small
\textbf{3. Etapa intermedia / Pirólisis primaria} (\SI{200}{\degreeCelsius}–\SI{600}{\degreeCelsius})
\begin{itemize}
  \item Se generan: \textbf{char primario}, gases condensables y no condensables.
  \item Se producen vapores precursores del bioaceite.
\end{itemize}

\vspace{0.2cm}
\textbf{4. Etapa final / Craqueo secundario} (\SI{300}{\degreeCelsius}–\SI{900}{\degreeCelsius})
\begin{itemize}
  \item Vapores condensables se craquean → más gases y \textbf{char secundario}.
  \item Si se eliminan rápidamente → se condensan como \textbf{bioaceite o alquitrán}.
  \item Temperaturas altas (\SI{>600}{\degreeCelsius}) favorecen producción de \ce{H2}.
  \item La \textbf{reacción shift} aumenta aún más el \ce{H2} por encima de \SI{900}{\degreeCelsius}.
\end{itemize}
\end{frame}


\begin{frame}{Observaciones Relevantes del Proceso de Pirólisis}
\small
\begin{itemize}
  \item A \textbf{baja temperatura}: alta producción de \ce{CO2}.
  \item \textbf{Hidrocarburos} máximos a \SI{450}{\degreeCelsius}, disminuyen > \SI{500}{\degreeCelsius} → se genera más \ce{H2}.
  \item El \textbf{char caliente} cataliza craqueo primario y secundario.
  \item \textbf{Remoción rápida de vapores condensables} evita craqueo → mejora el rendimiento de líquidos.
  \item \textbf{Conductividad térmica baja} (0.05–0.1 W/m·K): en leños grandes puede haber:
  \begin{itemize}
    \item Combustión externa,
    \item Mientras el interior sigue en etapa de secado.
    \item Fenómeno visible en incendios forestales.
  \end{itemize}
\end{itemize}
\end{frame}

\begin{frame}{Resumen de Etapas de la Pirólisis}
\small
\begin{center}
\begin{tikzpicture}[xscale=0.024, yscale=1, thick]
  % Ejes
  \draw[->] (90,0) -- (950,0) node[below] {Temperatura (\SI{500}{\celsius})};
  \draw[->] (90,0) -- (90,4) node[left] {Proceso};

  % Etapas
  \draw[fill=blue!20]  (100,0.5) rectangle (300,1.5) node[pos=.5] {\scriptsize Secado};
  \draw[fill=green!20] (100,1.6) rectangle (300,2.5) node[pos=.5] {\scriptsize Etapa Inicial (Torrefacción)};
  \draw[fill=yellow!30] (200,2.6) rectangle (600,3.5) node[pos=.5] {\scriptsize Pirólisis Primaria};
  \draw[fill=red!20]   (300,0.5) rectangle (900,1.5) node[pos=.5] {\scriptsize Craqueo Secundario};

  % Líneas verticales y anotaciones
  \foreach \x in {100,200,300,450,600,900} {
    \draw[dashed] (\x,0) -- (\x,3.8);
  }

  % Productos por etapa
  \node at (200,0.2) {\scriptsize H\textsubscript{2}O};
  \node at (250,1.3) {\scriptsize CO, CO\textsubscript{2}};
  \node at (400,3.7) {\scriptsize Bioaceite};
  \node at (700,1.3) {\scriptsize H\textsubscript{2}, Char};

\end{tikzpicture}
\end{center}
\smallskip
\textbf{Nota:} Las etapas se superponen parcialmente. La remoción rápida de vapores mejora el rendimiento líquido.
\end{frame}



\begin{frame}{Aspectos Químicos de la Pirólisis}
\small
\begin{itemize}
  \item La biomasa lignocelulósica tiene tres componentes principales:
  \begin{enumerate}
    \item Celulosa
    \item Hemicelulosa
    \item Lignina
  \end{enumerate}
  \item Cada uno presenta distintos comportamientos térmicos y químicos.
  \item El modelo de \textbf{Broido–Shafizadeh} (1979) explica la pirólisis de la celulosa mediante:
  \begin{itemize}
    \item Una \textbf{reacción intermedia} (I) → celulosa activa.
    \item Dos reacciones principales:
      \begin{itemize}
        \item II: \textit{Deshidratación} → char y gases (\ce{CO}, \ce{CO2}, \ce{H2O}).
        \item III: \textit{Despolimerización} → levoglucosano, alquitranes.
      \end{itemize}
  \end{itemize}
\end{itemize}
\end{frame}

\begin{frame}{Pirólisis de la Celulosa según Broido–Shafizadeh}
\small
\begin{itemize}
  \item \textbf{Reacción I:} genera un intermedio líquido —\textit{celulosa activa}— no estable, visible en pirólisis ablativa.
  \item \textbf{Reacción II (deshidratación):}
  \begin{itemize}
    \item Dominante a $T<\SI{300}{\degreeCelsius}$ y calentamiento lento.
    \item Produce \textbf{char, \ce{CO2}, \ce{CO}, \ce{H2O}}.
  \end{itemize}
  \item \textbf{Reacción III (despolimerización):}
  \begin{itemize}
    \item Dominante a altas tasas de calentamiento y $T>\SI{300}{\degreeCelsius}$.
    \item Produce vapores, \textbf{levoglucosano}, alquitrán.
  \end{itemize}
  \item \textbf{Reacción IV:} craqueo secundario de vapores → \textbf{char secundario}, alquitrán y gases.
  \item \textbf{Salida rápida del vapor} → evita craqueo → mayor rendimiento de líquidos (bioaceite).
\end{itemize}
\end{frame}

\begin{frame}{Pirólisis de Hemicelulosa y Lignina}
\small
\textbf{Hemicelulosa:}
\begin{itemize}
  \item Inicia degradación entre \SI{130}{\degreeCelsius}–\SI{194}{\degreeCelsius}, con mayor parte > \SI{180}{\degreeCelsius}.
  \item Produce más gas y menos alquitrán y char que la celulosa.
  \item Compuestos acuosos similares al ácido pirolenoso.
  \item Más humedad combinada y punto de ablandamiento más bajo.
\end{itemize}

\vspace{0.2cm}
\textbf{Lignina:}
\begin{itemize}
  \item Degradación inicia en \SI{280}{\degreeCelsius}, máximo entre \SI{350}{\degreeCelsius}–\SI{450}{\degreeCelsius}.
  \item Productos: \textbf{55\% char}, 15\% alquitrán, 20\% fase acuosa, 12\% gases.
  \item Alquitrán rico en fenoles → base para resinas ecológicas.
  \item Difícil de deshidratar.
\end{itemize}
\end{frame}

\begin{frame}{Constantes de Velocidad para Reacciones de Pirólisis}
\small
\begin{table}[h]
\centering
\begin{tabular}{@{}llcc@{}}
\toprule
\textbf{Reacción} & \textbf{Descripción} & $A_i$ (s\textsuperscript{–1}) & $E_i$ (kJ/mol) \\ \midrule
I & Degradación inicial (celulosa activa) & $2.83\times10^{19}$ & 243.0 \\
II & Deshidratación (char + gases) & $1.31\times10^{10}$ & 153.0 \\
III & Despolimerización (alquitrán) & $3.16\times10^{14}$ & 198.0 \\
IV & Craqueo secundario (gases, char) & $4.28\times10^{6}$ & 107.5 \\
\bottomrule
\end{tabular}
\caption{Constantes cinéticas del modelo de Broido–Shafizadeh. Tasa de reacción: $\frac{dm}{dt} = A_i (V_i^3 - V_i)\, e^{-E_i/RT}$}
\end{table}
\end{frame}


\begin{frame}{Modelo de Broido–Shafizadeh para la Celulosa}
\begin{figure}
    \centering
    \includegraphics[width=0.9\linewidth]{Energias Alternativas/4-Pyrolysis/img/modelo_broido_shafizadeh.png}
    \caption{Modelo de Broido–Shafizadeh}
   % \label{fig:enter-label}
\end{figure}
\end{frame}


\begin{frame}{Modelos Cinéticos Globales de Pirólisis}
\small
\begin{itemize}
  \item Para optimizar el proceso y maximizar los productos deseados, es necesario comprender la \textbf{cinética de la pirólisis}.
  \item La estimación de constantes cinéticas es difícil, especialmente para \textbf{pirólisis rápida}, por su carácter no estacionario.
  \item Se han propuesto tres modelos principales %\cite{Blasi1993}:
  \begin{enumerate}
    \item \textbf{Reacción global en una etapa}: una única reacción global.
    \item \textbf{Una etapa, múltiples reacciones}: varias reacciones paralelas, útil para distribución de productos.
    \item \textbf{Dos etapas, semiglobal}: incluye reacciones primarias y secundarias en serie.
  \end{enumerate}
  \item Para diseño ingenieril, el enfoque \textbf{“caja negra”} puede ser útil como primera aproximación.
\end{itemize}
\end{frame}

\begin{frame}{Modelo Global de Reacción Única}
\small
\textbf{Reacción simplificada:} \texttt{Biomasa → volátiles + char} (sin considerar cenizas ni humedad)

\vspace{0.2cm}
\textbf{Ecuaciones del modelo:}
\begin{align*}
\frac{dm_b}{dt} &= -k(m_b - m_c) \\
X &= \frac{m_0 - m_b}{m_0 - m_c} \\
\frac{dX}{dt} &= -k(1 - X) \\
X(t) &= 1 - A\,e^{-kt} \\
k &= \frac{E}{RT}
\end{align*}

\vspace{0.2cm}
\textbf{Notas:}
\begin{itemize}
  \item $m_b$: masa de biomasa sin pirolizar, $m_c$: masa final de char, $m_0$: masa inicial.
  \item $X$: conversión fraccional.
  \item $k$: constante de velocidad, $E$: energía de activación, $A$: factor preexponencial.
  \item Válido especialmente para \textbf{pirólisis lenta} bajo condiciones controladas.
\end{itemize}
\end{frame}

\begin{frame}{Constantes Cinéticas para el Modelo Global de Una Etapa}
\small
\begin{table}[h]
\centering
\begin{tabular}{@{}lccc@{}}
\toprule
\textbf{Combustible} & \textbf{Temperatura (K)} & $E$ (kJ/mol) & $A$ (s\textsuperscript{–1}) \\ \midrule
Celulosa & 520–1270 & 166.4 & $3.93 \times 10^{11}$ \\
Hemicelulosa & 520–1270 & 123.7 & $1.45 \times 10^{9}$ \\
Lignina & 520–1270 & 141.3 & $1.23 \times 10^{8}$ \\
Madera & 321–720 & 125.4 & $1.03 \times 10^{8}$ \\
Cáscara de almendra & 730–880 & 95–121 & $1.83 \times 10^{6}$ \\
Aserrín de haya & 450–700 & 84 (T < 600 K) & $2.33 \times 10^{4}$ \\
\bottomrule
\end{tabular}
\caption{Constantes para pirólisis a tasa de calentamiento lenta %\cite{Lewellen1977, Min1977, Font1990, Barooah1976}.}
\end{table}
\end{frame}

\begin{frame}{Transferencia de Calor en la Pirólisis}
\small
\begin{itemize}
  \item A temperaturas moderadas (\SI{300}{\degreeCelsius}–\SI{400}{\degreeCelsius}), la transferencia de calor no limita la pirólisis.
  \item A temperaturas mayores, \textbf{la conducción y convección internas influyen en la cinética global}.
  \item Modos de transferencia de calor involucrados %\cite{Babu2004b}:
  \begin{enumerate}
    \item \textbf{Conducción} en el interior de la partícula.
    \item \textbf{Convección} en los poros de la partícula.
    \item \textbf{Radiación y convección} desde la superficie externa.
  \end{enumerate}
  \item En reactores comerciales, el calor se transfiere a través de:
  \begin{itemize}
    \item Pared del reactor (vacío),
    \item Gases calientes (lecho arrastrado),
    \item Sólidos portadores de calor (lechos fluidizados).
  \end{itemize}
\end{itemize}
\end{frame}

\begin{frame}{Tamaño de Partícula y Transferencia de Masa}
\small
\begin{itemize}
  \item La baja conductividad térmica de la biomasa ($\sim$0.1 W/m·K) limita el calentamiento interno.
  \item Partículas grandes → calentamiento lento interior → más reacciones secundarias → \textbf{menor rendimiento líquido}.
  \item Shen et al. (2009): disminución del rendimiento de aceite con tamaños entre 0.3–1.5 mm; sin efecto apreciable al aumentar a 3.5 mm.
  \item \textbf{Transferencia de masa:}
  \begin{itemize}
    \item Una corriente de gas sobre el combustible ayuda a remover rápidamente los productos.
    \item Esto \textbf{minimiza el craqueo, repolimerización y recondensación} %\cite{Sensoz2008}.
  \end{itemize}
\end{itemize}
\end{frame}

\begin{frame}{¿La Pirólisis es Autotérmica?}
\small
\begin{itemize}
  \item \textbf{No completamente}: la energía liberada no basta para cubrir:
  \begin{itemize}
    \item Calentamiento del alimento y medios inertes.
    \item Reacciones endotérmicas.
    \item Pérdidas de calor.
  \end{itemize}
  \item Solución común: \textbf{combustión de gases no condensables y char} para generar calor.
  \item \textbf{Reacciones relevantes:}
  \begin{itemize}
    \item II (deshidratación): exotérmica.
    \item III (despolimerización) y IV (craqueo): endotérmicas.
    \item Pirólisis de hemicelulosa y lignina: exotérmica.
    \item Celulosa: endotérmica a $<\SI{400}{\degreeCelsius}$, exotérmica a mayor temperatura.
  \end{itemize}
\end{itemize}
\end{frame}

\begin{frame}{Reacciones Exotérmicas a Alta Temperatura (1000 K)}
\small
\begin{align*}
\ce{CO + 3H2 &-> CH4 + H2O} & \Delta H = -226~\text{kJ/mol} \\
\ce{CO + 2H2 &-> CH3OH} & \Delta H = -105~\text{kJ/mol} \\
\ce{0.17C6H10O5 &-> Char + 1.85H2O} & \Delta H = -80~\text{kJ/mol} \\
\ce{CO + H2O &-> CO2 + H2} & \Delta H = -42~\text{kJ/mol}
\end{align*}

\vspace{0.3cm}
\begin{itemize}
  \item A baja temperatura y corto tiempo de residencia → predominan reacciones endotérmicas ($\sim$2225 kJ/kg).
  \item A alta temperatura → predominan reacciones secundarias exotérmicas ($\sim$20 kJ/kg).
  \item En diseño, se puede \textbf{negligir el calor de reacción}, pero se debe considerar:
  \begin{itemize}
    \item Calor para calentar alimentación y gases.
    \item Calor para evaporación de productos.
  \end{itemize}
\end{itemize}
\end{frame}

\begin{frame}{Tipos de Pirolizadores y Variables Operativas}
\small
\begin{itemize}
  \item La selección del tipo de pirolizador depende del producto objetivo (char, líquido o gas).
  \item Variables que afectan el rendimiento:
  \begin{itemize}
    \item \textbf{Temperatura}, \textbf{tasa de calentamiento}, \textbf{tiempo de residencia del gas}.
  \end{itemize}
\end{itemize}

\vspace{0.3cm}
\begin{table}[h]
\small
\centering
\begin{tabular}{@{}lccc@{}}
\toprule
\textbf{Maximizar} & $T$ & Calentamiento & Tiempo gas \\ \midrule
Char   & Baja  & Lento & Largo \\
Líquido & Baja ($\sim$\SI{500}{\degreeCelsius}) & Rápido & Corto \\
Gas    & Alta  & Lento & Largo \\
\bottomrule
\end{tabular}
\caption{Efecto de variables operativas en el rendimiento de productos %\cite{Bridgwater1999, Demirbas2001}.}
\end{table}
\end{frame}

\begin{frame}{Clasificación General de Pirolizadores}
\small
\textbf{Según contacto gas–sólido:}
\begin{itemize}
  \item \textbf{Lecho fijo o móvil}
  \item \textbf{Lecho fluidizado burbujeante (BFB)}
  \item \textbf{Lecho fluidizado circulante (CFB)}
  \item \textbf{Ultrarrápido}
  \item \textbf{Cono rotatorio}
  \item \textbf{Ablativo}
  \item \textbf{Vacío}
\end{itemize}

\vspace{0.3cm}
\textbf{Objetivos modernos:} maximizar líquidos o gases en procesos continuos.
\end{frame}

\begin{frame}{Lecho Fijo y Lecho Fluidizado Burbujeante}
\small
\textbf{Lecho fijo:}
\begin{itemize}
  \item Operación discontinua.
  \item Calentamiento externo o combustión limitada (ej. horno colmena).
  \item Larga residencia → \textbf{alta producción de char}.
  \item Puede usarse gas inerte para barrido.
\end{itemize}

\vspace{0.2cm}
\textbf{Lecho fluidizado burbujeante (BFB):}
\begin{itemize}
  \item Biomasa (2–6 mm) alimentada a lecho de arena caliente.
  \item Buen control térmico y mezcla.
  \item \textbf{Rendimiento líquido: 70–75\,\%}.
  \item Requiere separación de char para evitar craqueo secundario.
  \item Fácil de escalar.
\end{itemize}
\end{frame}

\begin{frame}{CFB y Pirolizador Ultrarrápido}
\small
\textbf{CFB (Circulating Fluidized Bed):}
\begin{itemize}
  \item Mayor velocidad del gas, mayor capacidad.
  \item Reflujo interno y reciclaje de sólidos.
  \item Combustión externa del char para calentar el lecho.
\end{itemize}

\vspace{0.2cm}
\textbf{Ultrarrápido (UWO/Ensyn):}
\begin{itemize}
  \item Tiempo de mezcla: 10–20 ms; residencia: 70–200 ms.
  \item Inyección de nitrógeno o arena caliente a alta velocidad.
  \item \textbf{Rendimiento líquido hasta 90\,\%} \cite{Hulet2005}.
  \item Enfriamiento inmediato del vapor → minimiza craqueo.
\end{itemize}
\end{frame}

\begin{frame}{Otros Diseños: Ablativo, Cono y Vacío}
\small
\textbf{Ablativo:}
\begin{itemize}
  \item Presión entre partícula y pared caliente (contacto directo).
  \item \textbf{Rendimiento líquido: 80\,\%}.
  \item No requiere gas portador.
\end{itemize}

\textbf{Cono rotatorio:}
\begin{itemize}
  \item Alimentación desde la base con sólidos calientes.
  \item Excelente mezcla y calentamiento rápido (5000 K/s).
  \item \textbf{Rendimiento líquido: 60–70\,\%}.
  \item Diseño complejo para escalado.
\end{itemize}

\textbf{Vacío:}
\begin{itemize}
  \item Platos circulares apilados a distintas temperaturas.
  \item Sin gas portador, vapor extraído al vacío.
  \item \textbf{Rendimiento líquido: 35–50\,\%}, con alto char.
\end{itemize}
\end{frame}


% Diapositiva: Diseño para producción de bioaceite
\begin{frame}{Diseño para Producción de Bioaceite}
\small
\begin{itemize}
  \item La pirólisis puede producir hasta un \textbf{70\% de bioaceite} en peso seco.
  \item Factores clave para maximizar la producción líquida:
  \begin{itemize}
    \item \textbf{Alta velocidad de calentamiento}
    \item \textbf{Temperatura intermedia óptima} (\SI{500}{\degreeCelsius})
    \item \textbf{Tiempo de residencia corto} (0.1–2.0 s)
  \end{itemize}
  \item Balance térmico general:
  \begin{equation*}
    \text{Calor aportado} = \text{Calor requerido para pirólisis} + \text{Pérdidas}
  \end{equation*}
  \item Pérdidas estimadas: \textbf{10\% del calor de entrada}.
\end{itemize}
\end{frame}

% Diapositiva: Propiedades del bioaceite
\begin{frame}{Propiedades del Bioaceite}
\small
\begin{itemize}
  \item Poder calorífico: \textbf{13–18 MJ/kg}
  \item Composición: azúcares solubles, lignina pirolítica, agua.
  \item Contenido de oxígeno: \textbf{hasta 50\%}
  \item Mayor densidad (\textbf{\~1.3}) y viscosidad creciente con el tiempo.
  \item No es miscible con diésel ni autoignitable.
  \item Requiere \textbf{refinación o hidrocraqueo} para su uso como combustible.
\end{itemize}
\end{frame}

% Diapositiva: Introducción al Biochar
\begin{frame}{Producción de Biochar}
\small
\begin{itemize}
  \item Producto sólido de la \textbf{pirólisis lenta a temperatura moderada} (\SI{400}{\degreeCelsius}).
  \item Máximo rendimiento estimado: \textbf{35\% en peso}.
  \item Reacción estequiométrica típica (Antal, 2003):
  \begin{equation*}
    \text{C}_6\text{H}_{10}\text{O}_5 \rightarrow 3.74\text{C} + 2.65\text{H}_2\text{O} + 1.17\text{CO}_2 + 1.08\text{CH}_4
  \end{equation*}
\end{itemize}
\end{frame}

% Diapositiva: Beneficios del Biochar (1)
\begin{frame}{Beneficios del Biochar (1)}
\small
\begin{itemize}
  \item \textbf{Captura de carbono} y reducción de gases de efecto invernadero.
  \item Emisiones de carbono negativas.
  \item Sustitución de combustibles fósiles.
  \item Reducción de pérdida de nutrientes en suelos.
  \item Aumento de la productividad de suelos marginales.
  \item Mejora la retención de agua y estructura del suelo.
\end{itemize}
\end{frame}

% Diapositiva: Beneficios del Biochar (2)
\begin{frame}{Beneficios del Biochar (2)}
\small
\begin{itemize}
  \item Aumenta la \textbf{capacidad de intercambio catiónico (CEC)}.
  \item Mejora la calidad del agua reduciendo escorrentías contaminadas.
  \item Contribuye a la remediación de suelos y desertificación.
  \item Alternativa al "tumba y quema" de residuos agrícolas.
  \item Disminuye emisiones de \textbf{N$_2$O y CH$_4$} del suelo.
  \item Genera compensaciones de carbono y mayor rentabilidad agrícola.
\end{itemize}
\end{frame}

% Diapositiva: Propiedades del Biochar
\begin{frame}{Propiedades del Biochar}
\small
\begin{itemize}
  \item Superficie específica (área BET).
  \item pH del carbón.
  \item Capacidad de intercambio catiónico (CEC).
  \item Rendimiento de carbono en el biochar.
  \item Estas propiedades dependen de:
  \begin{itemize}
    \item \textbf{Temperatura de pirólisis}
    \item \textbf{Tipo de biomasa}
    \item \textbf{Tiempo de residencia}
  \end{itemize}
\end{itemize}
\end{frame}

\begin{frame}{Símbolos y Nomenclatura}
\small
\begin{tabular}{ll}
\textbf{Símbolo} & \textbf{Descripción} \\
\midrule
$A$ & Factor preexponencial (s\textsuperscript{–1}) \\
$E$ & Energía de activación (J/mol) \\
$k$ & Constante de velocidad de reacción (s\textsuperscript{–1}) \\
$m_b$ & Masa de biomasa en el tiempo $t$ (kg) \\
$m_c$ & Masa del residuo de char (kg) \\
$m_0$ & Masa inicial de la biomasa (kg) \\
$R$ & Constante universal de los gases (J/mol\,K) \\
$T$ & Temperatura (K) \\
$T_\text{pyr}$ & Temperatura de pirólisis (K) \\
$t$ & Tiempo (s) \\
$t_\text{heating}$ & Tiempo de calentamiento (s) \\
$t_r$ & Tiempo característico de reacción (s) \\
$X$ & Cambio fraccional en la masa de biomasa \\
\end{tabular}
\end{frame}


\end{document}